\documentclass[12pt]{article}
\title{Math 250, Thursday, Week 3}
\author{DZB}
\date{\today}

\begin{document}
\maketitle

``Tips" on mathematical writing. 
\\

\begin{itemize}
\item Your assignments should be written in prose. 
\item It should written in complete sentences; each sentence starts with a capitalized word (and in particular, doesn't start with a math symbol) and has punctuation. Write in paragraphs. 
\item On the homework we will deduct some points for sloppy prose. 
\end{itemize}

The intent with negations is that you `distrubute' the negation `all the way through'. 
\\

Technically, you can negate any sentence by adding "it is not true that" to the beginning. 
\\

All triangles are equilateral
\\

some triangles are not equilateral.\\

`This sentence has five words" ``It is not true that this sentence has five words"
\newpage

Proof by contradiction. Here is a template. 
\\

\begin{itemize}
\item We have a statement $P$ that we want to prove. 
\item There are `2 cases': either $P$ is true or $P$ is false. 
\item If we can rule out the case when $P$ is false, then $P$ is true. 
\item To begin, we assume that $P$ is false. (I.e., we assume $\neg P$.)
\item `Argue'; i.e., we exhibit some (correct) chain of reasoning, and end up with some statement $Q$.
\item (I.e., we write out a proof of $\neg P \Rightarrow Q$.)
\item We observe (or give a proof that) $Q$ is actually false. 
\item we conclude that $\neg P$ is false. 
\item Thus $P$ is true. 

\end{itemize}

How do you know when to proceed by contradiction? Ise contradiction when the negation `gives you something to work with'.
\\

Hieroglyph: $((\neg P \Rightarrow Q) \wedge \neg Q) \Rightarrow P$ (this is a statement form, and this is something that you can prove to be true by writing out a truth table).


\newpage

Definition. We say that a number $x$ is rational if there exist integers $a,b$ such that $b \neq 0$ and $x = a/b$. We write $\mathbf{Q}$ for rational numbers. We say that $x$ is irrational if $x \not \in \mathbf{Q}$.
\\

Examples: $1/3 \in \mathbf{Q}$; $\pi \not \in \mathbf{Q}$; $e$ is not rational; $4/6$ is rational; $4/6 = 2/3$. We say that $a/b$ is reduced if $a$ and $b$ have no common prime factors, and $a \geq 0$. (Example: $-3/5 = 3/(-5)$.) 
\\ 

Proposition: $\sqrt{2}$ is irrational.
\\

Sidebar: what does it mean to say that $\sqrt{2} \not \in \mathbf{Q}$? I.e., what is the negation of the definition of rational? For all $a,b \in \mathbf{Z}$ such that $b \neq 0$, $x \neq a/b$.
\\

We want to show that $\sqrt{2} = a/b$ has no solutions (with $a,b \in \mathbf{Z}$).

\newpage

Infinitude of primes. 
\\

Definition: a positive integer $p> 1$ is prime if the only postiive divisors are $1$ and $p$. A positive integer $n > 2$ is composite if it is not prime, i.e., there exists $d$ such $d \mid n$ and $1 < d < n$. Equivalently, $n$ is composite if we can write $n = ab$, where $a$ and $b$ are positive integers, and neither is equal to 1 or $n$.
\\

Theorem (Euclid): there are infinitely many primes. 
\\

Proof. 


\end{document}
